
\subsubsection{2.1 Execution-Feedback Reinforcement Learning for Code}

Reinforcement learning with execution feedback has been applied to code generation in several prior works. CodeRL \citep{Leetal2022} trained a code generation model on the APPS dataset using RL with execution outcomes as reward, demonstrating pass@1 improvements over supervised fine-tuning. CodeRL uses uniform random sampling from APPS and does not analyze how reward sparsity varies with problem difficulty or how it affects training dynamics.

RLEF \citep{Gehringetal2024} extends execution-feedback RL to function-level benchmarks including HumanEval, MBPP, and APPS, and shows that iterative refinement with execution-grounded reward outperforms both SFT and standard RLHF. Per-step training diagnostics such as advantage variance or reward density are not reported, and training data difficulty composition is not a research variable.

DeepSeek-R1 [DeepSeek-AI, 2025] applies GRPO with competitive programming data at large model scales (>30B parameters), achieving strong reasoning performance. At these scales, non-zero solve rates are more common, and the paper does not characterize the regime where the initializing policy cannot solve training problems.

None of these works measure per-step advantage variance or reward density as a function of task difficulty. This paper addresses that gap.

\subsubsection{2.2 Curriculum Learning}

Bengio et al. [2009] established that training on easier examples first can accelerate convergence and improve generalization. The mechanism is that easy examples provide stable gradient signal early in training, when the model is most sensitive to update direction. Subsequent work applied curriculum learning to neural machine translation \citep{Plataniosetal2019} and question answering [Sachan and Xing, 2016].

Curriculum learning in the context of execution-feedback GRPO has not been systematically studied. This paper identifies a mechanistic precondition for curriculum GRPO to function: the initializing policy must have a non-zero solve rate on easy-tier problems, so that early curriculum phases produce non-degenerate advantage groups. Whether this precondition holds for a given model and difficulty tier is an empirical question that is not answered in this paper.

Kumar et al. [2010] introduce self-paced learning, which adapts difficulty selection to the model's current performance rather than following a pre-defined schedule. This direction is complementary to the fixed-schedule approach considered here.

\subsubsection{2.3 GRPO}

GRPO \citep{Shaoetal2024} replaces the learned value baseline of PPO \citep{Schulmanetal2017} with group-relative advantage normalization. For each prompt, G completions are sampled; advantages are computed as (r_i − mean(r_group)) / std(r_group). This eliminates the need for a separate critic model but creates a structural dependency on reward variance within each group: when all completions receive the same reward, the gradient contribution is zero.

AlphaCode \citep{Lietal2022} demonstrates that model performance on competitive programming degrades sharply with problem difficulty tier, confirming that solve rate distributions are highly skewed for models of non-frontier scale. This paper quantifies how that skew translates into advantage variance collapse in GRPO.

\subsubsection{2.4 Positioning}

This work contributes a controlled measurement of GRPO advantage variance as a function of task granularity and reward structure, using training diagnostics (advantage variance, reward density) rather than only end-to-end benchmark outcomes. This diagnostic perspective is a prerequisite for principled design of curriculum approaches to execution-feedback GRPO.
